% !PRINTED 16
\weberpage{16}

%% [running head, not body text -- commented out 2026-08-07] Page furniture (running
%% head, folio, signature) is excluded from this edition by standing policy.
%% 16 \hfill \gesperrt{Erster Abschnitt.} \hfill \gesperrt{\S.\,5.}

%<<CONT
den isomorphen Gruppen (\S.\,1), nur mit dem Unterschiede, dass jedem Elemente $A$ nicht bloss ein Element $a$, sondern mehrere entsprechen. Diese Thatsache giebt Anlass, den Begriff des Isomorphismus zu erweitern, wie es durch folgende Definition geschieht:

\gesperrt{Man nennt eine Gruppe $P$ mit den Elementen $a, b, c \dots$ (mehrstufig) isomorph mit einer Gruppe $Q$ mit den Elementen $A, B, C \dots$, wenn beide Gruppen so auf einander bezogen werden können, dass jedem der Elemente $A, B, C \dots$ ein oder mehrere der Elemente $a, b, c \dots$ entsprechen, und zwar so, dass jedes der Elemente von $Q$ einem und nur einem der Elemente von $P$ entspricht, und dass, wenn $a$ und $A$, $b$ und $B$ einander entsprechen, $ab$ dem Elemente $AB$ entspricht.}
% [sic: Weber prints ``Elemente von Q''; the uniqueness condition requires P.]

So lässt sich zunächst beweisen, dass jedem der Elemente $A, B, C \dots$ eine gleiche Zahl von Elementen $a, b, c \dots$ entspricht und dass also \gesperrt{der Grad von $Q$ ein Theiler des Grades von $P$ ist}. Denn es sei $A$ das Einheitselement in $Q$, dem die Elemente $a, a_1 \dots a_{\nu-1}$ in $P$ entsprechen mögen. Diese letzteren Elemente müssen dann eine Gruppe vom Grade $\nu$ bilden, die wir mit $N$ bezeichnen wollen, weil nach der Definition des Isomorphismus das Element $a a_1$ dem Elemente $AA = A$ entsprechen muss. Ist dann $b$ ein dem Elemente $B$ entsprechendes Element von $P$, so müssen wiederum alle Elemente $Nb$ demselben Elemente $B$ aus $Q$ entsprechen. Denn jedes $ab$ muss dem Elemente $AB = B$ entsprechen. Sind $b_1, b$ zwei dem Elemente $B$ entsprechende Elemente, so entspricht $b_1 b^{-1} = a$ dem Elemente $A$, und ist also in $N$ enthalten, also ist $b_1 = ab$ in $Nb$ enthalten. Durch $Nb$ sind also alle Elemente, die dem Elemente $B$ entsprechen, erschöpft, und jedem Elemente von $Q$ entsprechen $\nu$ Elemente von $P$. Der Isomorphismus heisst \gesperrt{$\nu$-stufig}. Jedes Element $b^{-1}ab$ entspricht aber gleichfalls dem Einheitselemente $A$ und also ist $b^{-1}Nb = N$, d.\,h. $N$ ist Normaltheiler von $P$, und $Q$ ist einstufig isomorph mit $P/N$.

Wir sehen also, dass es einen anderen mehrstufigen Isomorphismus als den zwischen der complementären Gruppe eines Normaltheilers und der Gesammtgruppe nicht giebt. Man könnte aber den Begriff des Isomorphismus noch dahin erweitern, dass man zwei Gruppen, die zu einer dritten Gruppe $\mu$- und $\nu$-stufig
%CONT>>
